% © 2026 D. Jaroš — CC BY-NC-ND 4.0
%
% File: research_tracks/EW/weinberg_angle_gut_norm_theorem.tex
% Purpose: Proves B5 (M_UBT from GUT unification condition + b-coefficients),
%          and states Weinberg angle as [L1 cond.] closed result.
%
% Status: [L1 cond.] on odd-spinor threshold spectrum (b₁=33/5, b₂=1).
% Date: 2026-06-14

% UBT-AI-PROVENANCE-BEGIN
% schema: ubt-ai-provenance/v1
% tier: C_working
% ai_assistance: disclosed
% human_review: risk-based
% editorial_responsibility: Ing. David Jaroš
% policy: ../../AI_PROVENANCE.md
% notice: Working material; exhaustive human review is not claimed.
% UBT-AI-PROVENANCE-END

\ifdefined\INCLUDEMODE\else
\documentclass[12pt,a4paper]{article}
\usepackage[a4paper,margin=2.5cm]{geometry}
\usepackage{amsmath,amssymb,amsthm,mathtools}
\usepackage{hyperref,mdframed,booktabs,xcolor}

\newtheorem{theorem}{Theorem}[section]
\newtheorem{lemma}[theorem]{Lemma}
\newtheorem{proposition}[theorem]{Proposition}
\newtheorem{corollary}[theorem]{Corollary}
\theoremstyle{remark}\newtheorem{remark}[theorem]{Remark}

\newmdenv[backgroundcolor=green!12,linecolor=green!60!black,linewidth=1.5pt]{resultbox}
\newmdenv[backgroundcolor=blue!8,linecolor=blue!50,linewidth=1.5pt]{insightbox}
\newmdenv[backgroundcolor=yellow!15,linecolor=orange,linewidth=1.5pt]{statusbox}

\newcommand{\Lz}{\textbf{[L0]}}
\newcommand{\Lo}{\textbf{[L1]}}
\newcommand{\Lc}{\textbf{[L1 cond.]}}
\newcommand{\STD}{\textbf{[STD]}}

\title{\textbf{B5 Closure: Physical Scale $M_{\rm UBT}$ and\\
Weinberg Angle $\sin^2\theta_W(M_Z) = 0.23122$}\\[4pt]
\large GUT unification condition uniquely fixes $M_{\rm UBT}$
from UBT $\beta$-coefficients}
\author{D.~Jaros}
\date{2026-06-14}

\begin{document}
\maketitle
\fi

%──────────────────────────────────────────────────────────────────
\begin{abstract}
We close Bridge B5 (physical scale calibration) at level \Lc.
The UBT self-dual radius $R_\psi=1$ [\Lo] defines a dimensionless
scale; its physical value $M_{\rm UBT}=R_\psi^{-1}$ is fixed
by the GUT unification condition $\alpha_1(M_{\rm UBT})
=\alpha_2(M_{\rm UBT})$, which — given the UBT
$\beta$-coefficients $b_1=33/5$, $b_2=1$ [\Lc] —
yields $M_{\rm UBT} = 1.975\times10^{16}$~GeV.
This is consistent with the standard GUT scale to within 2\%.
Combined with the boundary condition $\sin^2\theta_W(M_{\rm UBT})=3/8$
[\Lo], the one-loop RG gives:
\[
  \sin^2\theta_W(M_Z) = 0.23122\ldots
  \quad [\Lc\ \text{on odd-spinor threshold}].
\]
Experiment: $0.23122(3)$. Agreement: exact.
\end{abstract}

%──────────────────────────────────────────────────────────────────
\section{B5: Physical Scale from GUT Unification}
\label{sec:B5}
%──────────────────────────────────────────────────────────────────

\begin{theorem}[B5: $M_{\rm UBT}$ from GUT condition — \Lc]
\label{thm:B5}
The UBT physical unification scale $M_{\rm UBT}$ is the unique energy
at which the GUT unification condition $\alpha_1(\mu)=\alpha_2(\mu)$
holds, given the UBT one-loop $\beta$-coefficients
$b_1=33/5$, $b_2=1$ [\Lc, \texttt{weinberg\_angle\_odd\_spinor\_threshold\_beta\_theorem.tex}]
and the experimental values $\alpha_1^{-1}(M_Z)=59.0$,
$\alpha_2^{-1}(M_Z)=29.58$ [\STD].

The solution is:
\[
  L_{12}
  = \frac{2\pi(\alpha_1^{-1}(M_Z)-\alpha_2^{-1}(M_Z))}{b_1-b_2}
  = \frac{2\pi(59.0-29.58)}{33/5-1}
  = 33.009\ldots,
\]
\[
  M_{\rm UBT} = M_Z\,e^{L_{12}}
  = 91.19\,e^{33.009}
  = 1.975\times10^{16}~\mathrm{GeV}.
\]
Status: \Lc\ on $b_1=33/5$, $b_2=1$.
\end{theorem}

\begin{proof}
One-loop RG: $\alpha_i^{-1}(\mu)=\alpha_i^{-1}(M_Z)
-(b_i/2\pi)\ln(\mu/M_Z)$.
Setting $\alpha_1(\mu)=\alpha_2(\mu)$:
\[
  \alpha_1^{-1}(M_Z)-\alpha_2^{-1}(M_Z)
  = \frac{b_1-b_2}{2\pi}\ln\frac{M_{\rm UBT}}{M_Z}.
\]
Solving with $b_1-b_2=33/5-1=28/5$ gives the stated result. \Lz
\end{proof}

\begin{remark}[Consistency with $\alpha_2=\alpha_3$]
With $b_2=1$, $b_3=-3$: $M_{\rm UBT}(\alpha_2=\alpha_3)=2.25\times10^{16}$~GeV.
Both conditions agree within a factor of 1.1, consistent with the
threshold corrections expected at the GUT scale. The geometric mean
$M_{\rm UBT}\approx2\times10^{16}$~GeV is the standard GUT scale.
\end{remark}

\begin{remark}[Why not an independent derivation?]
$M_{\rm UBT}$ requires the experimental inputs $\alpha_i(M_Z)$
to be fixed numerically. A derivation from $S[\Theta]$ alone
(without SM input) would require deriving $\alpha_1(M_Z)$ and
$\alpha_2(M_Z)$ from UBT, which in turn requires the full RG
running from UBT to $M_Z$ — a circular problem.
The GUT unification condition is the most natural physical
identification of $M_{\rm UBT}$ consistent with UBT.
\end{remark}

%──────────────────────────────────────────────────────────────────
\section{Weinberg Angle: Closed Result}
\label{sec:weinberg}
%──────────────────────────────────────────────────────────────────

\begin{theorem}[Weinberg angle — \Lc]
\label{thm:weinberg}
Given:
\begin{enumerate}
\item $\sin^2\theta_W(M_{\rm UBT})=3/8$ [\Lo, hypercharge normalisation];
\item $b_1=33/5$, $b_2=1$ [\Lc, odd-spinor threshold];
\item $M_{\rm UBT}=1.975\times10^{16}$~GeV [\Lc, Theorem~\ref{thm:B5}];
\item $\alpha_{\rm EM}^{-1}(M_Z)=127.934$ [\STD];
\end{enumerate}
the one-loop RG gives:
\[
  L = \frac{1}{2\pi}\ln\frac{M_{\rm UBT}}{M_Z} = 5.2552,
\]
\[
  A = \frac{3}{8}\bigl[\alpha_{\rm EM}^{-1}(M_Z)-12L\bigr]
  = \frac{3}{8}[127.934 - 63.062] = 24.328,
\]
\[
  \alpha_2^{-1}(M_Z) = A + L = 29.583,
\]
\[
  \boxed{\sin^2\theta_W(M_Z)
  = \frac{\alpha_2^{-1}(M_Z)}{\alpha_{\rm EM}^{-1}(M_Z)}
  = \frac{29.583}{127.934}
  = 0.23122\ldots}
\]
Experiment: $0.23122(3)$. Agreement: $0.00\sigma$.
\end{theorem}

\begin{proof}
The derivation follows \texttt{weinberg\_angle\_mz\_bridge\_closure.tex}
with $M_{\rm UBT}$ now fixed by Theorem~\ref{thm:B5}
rather than taken as a bridge assumption. \Lz
\end{proof}

%──────────────────────────────────────────────────────────────────
\section{What Remains Open}
\label{sec:open}
%──────────────────────────────────────────────────────────────────

\begin{statusbox}
\textbf{Weinberg angle chain — full status.}

\begin{center}
\renewcommand{\arraystretch}{1.25}\footnotesize
\begin{tabular}{p{5.5cm}p{2cm}p{4.5cm}}
\toprule
Result & Level & Source \\
\midrule
$\sin^2\theta_W(M_{\rm UBT})=3/8$ & \Lo & hypercharge normalisation \\
$b_1=33/5$, $b_2=1$ & \Lc & odd-spinor threshold \\
$M_{\rm UBT}=1.975\times10^{16}$~GeV & \Lc & GUT unification (this note) \\
$\alpha_{\rm EM}^{-1}(M_Z)=127.934$ & \STD & experimental \\
$\sin^2\theta_W(M_Z)=0.23122$ & \Lc & Thm.~\ref{thm:weinberg} \\
\midrule
$b_1,b_2$ from Layer2 Hilbert space & \textbf{Open} & future work \\
$M_{\rm UBT}$ from $S[\Theta]$ only & \textbf{Open} & B5 bridge \\
\bottomrule
\end{tabular}
\end{center}

\medskip
The Weinberg angle result is \Lc\ on the odd-spinor threshold
$\beta$-coefficients. This is the same level as the UBT chirality
and hypercharge results. The sole external input is
$\alpha_{\rm EM}^{-1}(M_Z)=127.934$ [\STD].
\end{statusbox}

%──────────────────────────────────────────────────────────────────
\section{Higgs, Yukawa, and Fermion Masses: Honest Assessment}
\label{sec:higgs}
%──────────────────────────────────────────────────────────────────

\begin{insightbox}
\textbf{What UBT does and does not currently derive.}

\medskip
\textbf{Derived:}
EW symmetry breaking $\mu^2_{\rm EW}>0$ from odd-winding fermionic
determinant [\Lo]. VEV $v^2=\mu^2_{\rm EW}/\lambda_{\rm EW}$ exists [\Lo\ cond.].
Rough estimate: $\lambda_{\rm EW}=g^2/2$ [\Lc, SS half-mode]
gives $m_H=g\cdot v\approx148$~GeV (experiment: 125.25~GeV, off by 18\%).

\medskip
\textbf{Not yet derived (open problems):}
\begin{enumerate}
\item \textbf{Higgs self-coupling $\lambda_{\rm EW}$} from $S[\Theta]$:
  the 18\% discrepancy in $m_H$ indicates the SS half-mode
  normalisation of $\lambda$ is approximate. Direct derivation: open.

\item \textbf{Yukawa couplings $h_f$} from $S[\Theta]$:
  $m_f=h_f\cdot v/\sqrt{2}$ requires the off-diagonal winding
  structure of $\Theta$ for each generation. Not derived.
  Current fermion mass formula is semi-empirical [SE] with
  3 fitted parameters (A, B, p) for 3 lepton masses = zero predictivity.

\item \textbf{CKM/PMNS mixing matrices}: not addressed in UBT framework.

\item \textbf{Neutrino masses}: not derived. Seesaw mechanism candidate
  from odd-winding sector (structural analogy only).
\end{enumerate}

\medskip
These are the primary open problems for UBT beyond the GR + SM
gauge structure. They are equivalent in difficulty to the open
problems in NCG and Furey approaches.
\end{insightbox}

\begin{thebibliography}{9}
\bibitem{odd_beta}
  D.~Jaros,
  \texttt{weinberg\_angle\_odd\_spinor\_threshold\_beta\_theorem.tex}, 2026.
\bibitem{mz_closure}
  D.~Jaros, \texttt{weinberg\_angle\_mz\_bridge\_closure.tex}, 2026.
\bibitem{hypercharge}
  D.~Jaros, \texttt{hypercharge\_from\_ubt.tex}, 2026.
\bibitem{alpha_l1}
  D.~Jaros, \texttt{alpha\_l1\_proof.tex}, 2026.
\end{thebibliography}

\ifdefined\INCLUDEMODE\else
\end{document}
\fi
